\documentclass[10pt,a4paper]{amsart}
\usepackage{amsmath,amssymb,amsthm,graphicx}
\usepackage{verbatim}
\newenvironment{reference}{\comment}{\endcomment}
%\newenvironment{reference}{}{}
\newenvironment{slogan}{\comment}{\endcomment}
\newenvironment{history}{\comment}{\endcomment}

% For commutative diagrams we use Xy-pic
\usepackage[all]{xy}

% We use 2cell for 2-commutative diagrams.
\xyoption{2cell}
% Theorem environments.
%
\theoremstyle{plain}
\newtheorem{theorem}[subsection]{Theorem}
\newtheorem{proposition}[subsection]{Proposition}
\newtheorem{lemma}[subsection]{Lemma}

\theoremstyle{definition}
\newtheorem{definition}[subsection]{Definition}
\newtheorem{example}[subsection]{Example}
\newtheorem{exercise}[subsection]{Exercise}
\newtheorem{situation}[subsection]{Situation}

\theoremstyle{remark}
\newtheorem{remark}[subsection]{Remark}
\newtheorem{remarks}[subsection]{Remarks}

\numberwithin{equation}{subsection}

% Macros
%
\def\lim{\mathop{\mathrm{lim}}\nolimits}
\def\colim{\mathop{\mathrm{colim}}\nolimits}
\def\Spec{\mathop{\mathrm{Spec}}}
\def\Hom{\mathop{\mathrm{Hom}}\nolimits}
\def\Ext{\mathop{\mathrm{Ext}}\nolimits}
\def\SheafHom{\mathop{\mathcal{H}\!\mathit{om}}\nolimits}
\def\SheafExt{\mathop{\mathcal{E}\!\mathit{xt}}\nolimits}
\def\Sch{\mathit{Sch}}
\def\Mor{\mathop{\mathrm{Mor}}\nolimits}
\def\Ob{\mathop{\mathrm{Ob}}\nolimits}
\def\Sh{\mathop{\mathit{Sh}}\nolimits}
\def\NL{\mathop{N\!L}\nolimits}
\def\CH{\mathop{\mathrm{CH}}\nolimits}
\def\proetale{{pro\text{-}\acute{e}tale}}
\def\etale{{\acute{e}tale}}
\def\QCoh{\mathit{QCoh}}
\def\Ker{\mathop{\mathrm{Ker}}}
\def\Im{\mathop{\mathrm{Im}}}
\def\Coker{\mathop{\mathrm{Coker}}}
\def\Coim{\mathop{\mathrm{Coim}}}

% Boxtimes
%
\DeclareMathSymbol{\boxtimes}{\mathbin}{AMSa}{"02}

%
% Macros for moduli stacks/spaces
%
\def\QCohstack{\mathcal{QC}\!\mathit{oh}}
\def\Cohstack{\mathcal{C}\!\mathit{oh}}
\def\Spacesstack{\mathcal{S}\!\mathit{paces}}
\def\Quotfunctor{\mathrm{Quot}}
\def\Hilbfunctor{\mathrm{Hilb}}
\def\Curvesstack{\mathcal{C}\!\mathit{urves}}
\def\Polarizedstack{\mathcal{P}\!\mathit{olarized}}
\def\Complexesstack{\mathcal{C}\!\mathit{omplexes}}
% \Pic is the operator that assigns to X its picard group, usage \Pic(X)
% \Picardstack_{X/B} denotes the Picard stack of X over B
% \Picardfunctor_{X/B} denotes the Picard functor of X over B
\def\Pic{\mathop{\mathrm{Pic}}\nolimits}
\def\Picardstack{\mathcal{P}\!\mathit{ic}}
\def\Picardfunctor{\mathrm{Pic}}
\def\Deformationcategory{\mathcal{D}\!\mathit{ef}}
\usepackage{cleveref}
\crefname{originallemma}{original lemma}{original lemmas}
\Crefname{originallemma}{Original lemma}{Original lemmas}
\makeatletter
\newlabel{lemma-branches}{{0DRB}{source}{}{sourcepointer-1}{}}
\newlabel{lemma-branches@cref}{{[originallemma][0][]0DRB}{[1][0][]source}{}{}{}}
\makeatother
\crefname{originallemma}{original lemma}{original lemmas}
\Crefname{originallemma}{Original lemma}{Original lemmas}
\makeatletter
\newlabel{lemma-Artin-approximation-by-smooth-morphism}{{0DR0}{source}{}{sourcepointer-2}{}}
\newlabel{lemma-Artin-approximation-by-smooth-morphism@cref}{{[originallemma][0][]0DR0}{[1][0][]source}{}{}{}}
\makeatother
\crefname{originalreference}{original reference}{original references}
\Crefname{originalreference}{Original reference}{Original references}
\makeatletter
\newlabel{more-algebra-proposition-finite-type-over-G-ring}{{07PV}{source}{}{sourcepointer-3}{}}
\newlabel{more-algebra-proposition-finite-type-over-G-ring@cref}{{[originalreference][0][]07PV}{[1][0][]source}{}{}{}}
\makeatother
\crefname{originalreference}{original reference}{original references}
\Crefname{originalreference}{Original reference}{Original references}
\makeatletter
\newlabel{algebra-lemma-pullback-module}{{02M1}{source}{}{sourcepointer-4}{}}
\newlabel{algebra-lemma-pullback-module@cref}{{[originalreference][0][]02M1}{[1][0][]source}{}{}{}}
\makeatother
\title{The formal-branch map preserves multiplicities when the base is a G-ring}
\author{The Stacks Project Authors}
\date{stacks-project:stacks-geometry.tex:local-snapshot-20261012; GNU FDL 1.2 or later}
\begin{document}
\maketitle
\noindent\textit{Editorial note.} All delivered mathematical text is the author's own, copied verbatim from the pinned Stacks Project chapter file stacks-geometry.tex; no statement or proof step was rewritten, reordered or removed.\par\medskip
\noindent\textit{Editorial note.} Span recomputation by label pairing (LANE\_BRIEF 10.2): in this chapter the author's proof environment follows the statement environment as a standalone block; the spec covers the complete environment pair (statement span, proof span) as recorded in the pool candidate.\par\medskip
\noindent\textit{Editorial note.} Occupation scan: combined claim\_key + statement\_anchor + generic-line scan over all registration files (active and withdrawn) found 0 hits for stacks-geometry.tex (ge-occupied.json, 2026-10-14); 207 route=geometry / chapter-list field hits in focused-seeds are alias noise (route name and alphabetical chapter enumerations), excluded per the ladder alias-blind-spot caveat.\par\medskip
\ifdefined\hypertarget\hypertarget{sourcepointer-1}{}\fi
\noindent\textit{Source reference.} Original lemma 0DRB: sources/stacks-chow/tags.txt, source lines 17216--17216. Cross-chapter standard prerequisite cited by the author inside the counted span; it is used as a pointer and its own proof is not part of this item.\par\medskip
\ifdefined\hypertarget\hypertarget{sourcepointer-2}{}\fi
\noindent\textit{Source reference.} Original lemma 0DR0: sources/stacks-chow/tags.txt, source lines 17205--17205. Cross-chapter standard prerequisite cited by the author inside the counted span; it is used as a pointer and its own proof is not part of this item.\par\medskip
\ifdefined\hypertarget\hypertarget{sourcepointer-3}{}\fi
\noindent\textit{Source reference.} Original reference 07PV: sources/stacks-chow/tags.txt, source lines 9810--9810. Cross-chapter standard prerequisite cited by the author inside the counted span; it is used as a pointer and its own proof is not part of this item.\par\medskip
\ifdefined\hypertarget\hypertarget{sourcepointer-4}{}\fi
\noindent\textit{Source reference.} Original reference 02M1: sources/stacks-chow/tags.txt, source lines 3428--3428. Cross-chapter standard prerequisite cited by the author inside the counted span; it is used as a pointer and its own proof is not part of this item.\par\medskip
\begin{lemma}
\label{lemma-branches-multiplicity}
Let $\mathcal{X}$ be an algebraic stack locally of finite type
over a locally Noetherian scheme $S$. Let $x_0 : \Spec(k) \to \mathcal{X}$
is a morphism where $k$ is a field of finite type over $S$ with
image $s \in S$. If $\mathcal{O}_{S, s}$ is a G-ring, then
the map of Lemma \ref{lemma-branches} preserves multiplicities.
\end{lemma}

\begin{proof}
By Lemma \ref{lemma-Artin-approximation-by-smooth-morphism} we
may assume there is a smooth morphism $U \to \mathcal{X}$
where $U$ is a scheme and a $k$-valued point $u_0$ of $U$
such that $\mathcal{O}_{U, u_0}^\wedge$ is a versal ring to
$\mathcal{X}$ at $x_0$. By construction of our map in
the proof of Lemma \ref{lemma-branches} (which simplifies
greatly because $A = \mathcal{O}_{U, u_0}^\wedge$) we find
that it suffices to show: the multiplicity of an irreducible
component of $U$ passing through $u_0$ is the same as the
multiplicity of any irreducible component of
$\Spec(\mathcal{O}_{U, u_0}^\wedge)$ mapping into it.

\medskip\noindent
Translated into commutative algebra we find the following:
Let $C = \mathcal{O}_{U, u_0}$. This is essentially of finite type
over $\mathcal{O}_{S, s}$ and hence is a G-ring (More on Algebra,
Proposition \ref{more-algebra-proposition-finite-type-over-G-ring}).
Then $A = C^\wedge$. Therefore $C \to A$ is a regular ring map.
Let $\mathfrak q \subset C$ be a minimal prime and let $\mathfrak p \subset A$
be a minimal prime lying over $\mathfrak q$. Then
$$
R = C_\mathfrak p \longrightarrow A_\mathfrak p = R'
$$
is a regular ring map of Artinian local rings. For such a ring map
it is always the case that
$$
\text{length}_R R = \text{length}_{R'} R'
$$
This is what we have to show because the left hand side
is the multiplicity of our component on $U$ and the right hand
side is the multiplicity of our component on $\Spec(A)$.
To see the equality, first we use that
$$
\text{length}_R(R) \text{length}_{R'}(R'/\mathfrak m_R R') =
\text{length}_{R'}(R')
$$
by Algebra, Lemma \ref{algebra-lemma-pullback-module}. Thus it suffices
to show $\mathfrak m_R R' = \mathfrak m_{R'}$, which is a consequence
of being a regular homomorphism of zero dimensional local rings.
\end{proof}

\end{document}
